\section{Introduction}
\label{sec:intro}

\subsection{Counting maps between graphs}

Let $G$ and $H$ be finite graphs, where $H$ may carry loops. A
\emph{homomorphism} $\varphi \colon G \to H$ is a map
$V(G) \to V(H)$ that sends edges to edges: if $uv \in E(G)$ then
$\varphi(u)\varphi(v) \in E(H)$. We write $\Hom(G,H)$ for the number of
such maps.

The quantity $\Hom(G,H)$ specialises to many familiar counts. If
$H = K_q$ is complete, it counts proper $q$-colorings; if $H$ is a
looped edge, it counts independent sets; if $H$ is a path, it counts
height functions with bounded increments. We call $H$ the
\emph{target}, and we often call $\Hom(G,H)$ the number of
\emph{$H$-colorings} of $G$.

Two extremal questions arise: among all trees $T$ on $n$ vertices,
which maximizes $\Hom(T,H)$, and which minimizes it? Sidorenko~\cite{sidorenko1991}
settled the first: the star $K_{1,n-1}$ is always a maximizer for every
target $H$. The second remains open in general.

\subsection{The path conjecture and its failure}

Galvin and Tetali~\cite{galvin-tetali} proposed the natural companion:
that the path $P_n$ should always be a minimizer. The intuition is
appealing. A homomorphism from a tree is built by propagating
constraints outward from a root, and the path is the tree that spreads
those constraints along the longest possible chain, giving each vertex
the fewest independent choices. Csikv\'ari and Lin~\cite{csikvari-lin}
developed the systematic tool for this circle of ideas, the
KC-transformation, and proved the conjecture for large classes of
targets.

The conjecture is false. Leontovich~\cite{leontovich1989} exhibited a
target for which some tree beats the path, and
Galvin, McMillon, Redlich and Nir~\cite{galvin2025} gave explicit
bipartite counterexamples. We call $H$ a \emph{Leontovich graph} when
the path fails to minimize, and we call the least such $n$ the
\emph{crossover threshold}. A companion paper~\cite{galvin2026}
introduces the dual class of \emph{Hoffman--London graphs}; here we ask
how small a Leontovich graph can be, and how early a crossover can
occur.

What makes the problem hard is that the failure is faint. In every known
example the counts differ by a relative margin of about $10^{-4}$ or
less, and the winner is usually just a single pendant move away from the
path. The signal is a near-cancellation of exponentially growing terms,
so crude estimates miss it and floating-point noise can erase it. Much
of what follows is therefore about comparisons that survive at that
scale, and about marking proofs separately from bounded searches.

\subsection{What this paper does}

We attack the problem in three ways.

\emph{Analytically.} We isolate an eigenvalue condition that reduces each
odd near-path comparison to a single base-case sign check
(Theorem~\ref{thm:spectral}), determine the exact odd crossover
threshold of the smallest previously known Leontovich tree
(Theorem~\ref{thm:threshold}), and show that the bipartite double cover
preserves both the Leontovich property and its asymptotic strengthening
(Theorem~\ref{thm:double_cover}).

\emph{Computationally.} We sweep large families of targets --- all
connected simple graphs on at most $11$ vertices, all connected
bipartite graphs on at most $15$ vertices, and all symmetric trees
below $78$ vertices --- against the near-path source trees defined
below, and we record the results as bounded sweeps rather than theorems.

\emph{By synthesis.} Where enumeration runs out, we encode the
crossover condition as an integer constraint problem and hand it to an
SMT solver. This produces the $18$-vertex depth-$2$ witness; the
$15$-vertex depth-dependent witness comes from exhaustive bipartite
enumeration.

\subsection{Standards of evidence}
\label{sec:evidence}

Because the arguments mix proof with computation, and because earlier
drafts contained claims that failed exact arithmetic, we label every
substantive assertion with one of four tags and use them consistently.

\begin{description}[leftmargin=2.4em,style=nextline,font=\normalfont\bfseries]
  \item[(P) Proved.] An unconditional proof, quantifiers included.
    Stated as \textbf{Theorem}, \textbf{Proposition}, \textbf{Lemma} or
    \textbf{Corollary}.
  \item[(E) Exact finite computation.] A numerical claim about named
    objects, verified in exact integer or certified
    arbitrary-precision arithmetic.
  \item[(S) Bounded sweep.] A negative result of the form ``no member
    of family $\mathcal{F}$ exhibits behaviour $B$ within the tested
    window''. This certifies the window, not the family; we never state it as a
    nonexistence theorem.
  \item[(N) High-precision numerics.] A floating-point computation with
    a residual-based error estimate but no outward-rounded interval or
    exact enclosure. This is numerical evidence, not a proof.
\end{description}

The distinction between (P) and (S) matters most. A near-path sweep
certifies only that the tested trees fail to beat the path; it does not
certify that no tree does. Accordingly, we reserve nonexistence claims
for proved results and state the scope of each sweep explicitly.

\subsection{Summary of contributions}

\begin{enumerate}[label=(\roman*)]
  \item \textbf{(E)} Verification of the path-minimizer theorem for the
    odd paths $H \in \{P_3,P_5,P_7,P_9\}$ over all source trees on
    $n \le 22$ vertices. For $P_3$ and $P_5$ the path is minimal but
    not unique; from $P_7$ onward it is uniquely minimal. The
    arithmetic mechanism is that, with $k$ the path-length parameter
    and $\lambda_1$ the first eigenvalue of the active quotient,
    $\lambda_1^2$ is rational exactly for
    $k \in \{1,2,3,5\}$. This resolves the question computationally
    through $n \le 22$ and identifies the arithmetic mechanism behind
    the observed transition in Problem~6.3 of~\cite{galvin2026}
    (Section~\ref{sec:odd-path}).
  \item \textbf{(P)} An eigenvalue criterion (Theorem~\ref{thm:spectral}):
    for a rooted bipartite target whose active quotient has a single
    positive eigenvalue, each fixed-depth odd near-path comparison has
    constant sign. For all $2$-orbit and $3$-orbit symmetric trees this
    reduces detection to one base-case check
    (Corollary~\ref{cor:2_and_3_orbit_trees}).
  \item \textbf{(P)} Among odd $n$, the crossover threshold of
    $T(7,1,9)$ is exactly $n = 13$, and it is permanent: the sign changes
    once and never again (Theorem~\ref{thm:threshold}). This threshold was already
    known in correspondence. The companion impossibility at $n = 5$ for
    the spherically symmetric family $T(x,y,z)$ is a polynomial
    identity (Proposition~\ref{prop:n5}).
  \item \textbf{(S)} All $1{,}018{,}680{,}329$ connected simple graphs on
    $m \le 11$ vertices, and all $608{,}930{,}559$ connected bipartite
    graphs on $m \le 15$ vertices, tested against the near-path family
    with $n \le 200$, $d \le 20$. No simple target below $12$ vertices
    produced a near-path crossover. This bounds, but does not settle,
    Problem~4.3 of~\cite{galvin2025} (Section~\ref{sec:bounds}).
  \item \textbf{(E)} Two small bipartite witnesses: a $15$-vertex graph
    $H^*$ whose earliest crossover is at $(n,d) = (49,16)$, and an
    $18$-vertex graph $H_{18}$ whose first near-path crossover is at
    $(n,d)=(15,4)$ with margin $\Delta=500{,}576$. It crosses against
    the classical depth-$2$ near-path at $n=17$ with margin
    $\Delta = 5{,}068{,}778$ (Section~\ref{sec:discussion}).
  \item \textbf{(P)} The bipartite double cover preserves both the
    Leontovich and the strongly Leontovich property
    (Theorem~\ref{thm:double_cover}). Applied to the looped tree
    $\hat T(1,35,1,50)$, this yields a \emph{simple} strongly Leontovich
    bipartite graph on $3{,}644$ vertices, whose first even crossover we
    locate exactly at $n = 17{,}340$ --- evidence that even-$n$
    crossovers do not require a loop, only enough room to delay the
    cancellation.
  \item \textbf{(N)} A perturbed double cover on $6{,}806$ vertices whose
    estimated leading-coefficient ratio is $0.999952140676 < 1$,
    a candidate simple non-bipartite strongly Leontovich target.
\end{enumerate}

Section~\ref{sec:prelim} fixes notation and records the two
transfer-matrix identities that every computation in the paper rests on.
